\begin{abstract}
Compiling GUI procedures that agents execute repeatedly into programs can reduce their token costs.
However, deciding when to compile poses two challenges.
First, the compilation cost is unknown before an attempt, which may fail after consuming tokens.
Second, future use is unknown because tasks may stop repeating or GUI changes may break the program.
To address these challenges, we propose PACE (Payback-Aware Compilation from Experience), a system with a measurement protocol and an online compilation protocol.
The measurement protocol measures compilation costs and per-use savings to estimate how many uses a program needs to pay back its compilation cost.
Using these measurements, the online protocol estimates whether future savings justify compilation based on past task arrivals and compilation outcomes.
It checks execution and compilation charges against a cumulative budget before allowing either action.
Under stated action-cost assumptions, total cost after each arrival is at most $1+\epsilon$ times the cost of running every task with the agent.
Based on measured costs, verified programs need an estimated $1$--$14$ uses to recover the cost of a successful compilation.
In experiments using recorded task arrivals, PACE reduces token costs by $21.9\%$ compared with ReAct, $26.9\%$ with AutoRPA, and $21.6\%$ with ToolPro on average ($\epsilon=0.25$).
\end{abstract}
